\documentclass[11pt,a4paper]{article}

\usepackage[top=2cm,bottom=2.2cm,left=2.2cm,right=2.2cm]{geometry}
\usepackage[utf8]{inputenc}
\usepackage[T1]{fontenc}
\usepackage{lmodern}
\usepackage{microtype}
\usepackage{amsmath,amssymb}
\usepackage{booktabs}
\usepackage{tabularx}
\usepackage{xcolor}
\usepackage{fancyhdr}
\usepackage{enumitem}
\usepackage{titlesec}
\usepackage{array}
\usepackage{float}
\usepackage[colorlinks=true,linkcolor=navyblue,urlcolor=navyblue,citecolor=navyblue]{hyperref}

% Definisi Warna Elegan
\definecolor{navyblue}{RGB}{16, 44, 87}
\definecolor{darkslate}{RGB}{30, 41, 59}
\definecolor{accentblue}{RGB}{41, 128, 185}
\definecolor{softblue}{RGB}{240, 245, 255}
\definecolor{graytext}{RGB}{100, 116, 139}

% Format Heading
\titleformat{\section}{\large\bfseries\color{navyblue}}{\thesection.}{0.6em}{}[\titlerule]
\titleformat{\subsection}{\normalsize\bfseries\color{darkslate}}{\thesubsection}{0.5em}{}
\titleformat{\subsubsection}{\small\bfseries\color{accentblue}}{\thesubsubsection}{0.4em}{}
\titlespacing*{\section}{0pt}{12pt}{4pt}
\titlespacing*{\subsection}{0pt}{8pt}{3pt}
\titlespacing*{\subsubsection}{0pt}{6pt}{2pt}

% Header dan Footer
\pagestyle{fancy}
\fancyhf{}
\fancyhead[L]{\small\color{graytext}Laporan Kemajuan Tugas Akhir --- M. Hakim Ghifari (6182201096)}
\fancyhead[R]{\small\color{graytext}Teknik Informatika FTIS UNPAR}
\fancyfoot[C]{\thepage}
\renewcommand{\headrulewidth}{0.4pt}
\renewcommand{\footrulewidth}{0pt}

\setlist{nosep,leftmargin=1.4em}

\begin{document}

% -----------------------------------------------------------------------------
% KOP LAPORAN RESMI
% -----------------------------------------------------------------------------
\begin{center}
    {\large\textbf{\color{navyblue}PROGRAM STUDI TEKNIK INFORMATIKA}}\\[2pt]
    {\normalsize\textbf{FAKULTAS SAINS --- UNIVERSITAS KATOLIK PARAHYANGAN}}\\[6pt]
    {\LARGE\bfseries\color{navyblue}LAPORAN KEMAJUAN PENELITIAN TUGAS AKHIR}\\[3pt]
    {\small\textit{Rekapitulasi Progres Eksperimen, Implementasi Metodologi, \& Penulisan Naskah Skripsi}}
\end{center}

\vspace{0.1cm}
\noindent\hrulefill

\vspace{0.2cm}
\noindent
\begin{tabularx}{\textwidth}{@{}llX@{}}
    \textbf{Topik Penelitian} & : & \textbf{Analisis Data Satelit Multi-Spektral Time Series untuk Mengidentifikasi Padi Rojolele Srinuk Berdasarkan Pola Pertumbuhannya di Klaten} \\
    \textbf{Mahasiswa}        & : & \textbf{Muhammad Hakim Ghifari Akiyat} (NPM: \textbf{6182201096}) \\
    \textbf{Dosen Pembimbing} & : & \textbf{Prof. Dr. Veronica Sri Moertini, Ir., M.T.} \\
    \textbf{Tanggal Laporan}  & : & \textbf{30 September 2026} \\
    \textbf{Status Progres}   & : & \textbf{Validasi Model Selesai (Akurasi 94,23\%--98,08\%) | WebGIS \& 110+ Grafik Siap | Bab 1--2 Selesai} \\
\end{tabularx}

\vspace{0.2cm}
\noindent\hrulefill

% -----------------------------------------------------------------------------
% 1. RINGKASAN EKSEKUTIF (DASHBOARD)
% -----------------------------------------------------------------------------
\section{Ringkasan Eksekutif Kemajuan Penelitian (Progress Dashboard)}

Laporan kemajuan ini mendokumentasikan pencapaian komprehensif penelitian Tugas Akhir mengenai deteksi dan identifikasi \textbf{Padi Varietas Unggul Lokal Rojolele Srinuk} di Kecamatan Delanggu, Kabupaten Klaten, Jawa Tengah. Penelitian ini mengadaptasi dan mengembangkan metodologi skripsi \textbf{Vico Pratama (2025)} dengan menerapkan fusi citra multi-sensor historis (1988--2026), segmentasi superpixel spasial \textit{Simple Non-Iterative Clustering} (SNIC), dan pemodelan klasifikasi deret waktu berjenjang \textit{K-Nearest Neighbors} berbasis \textit{Dynamic Time Warping} (KNN-DTW).

\begin{table}[htbp]
    \centering\small
    \caption{Dasbor Capaian Progres Komponen Kerja Penelitian (Per 30 September 2026)}
    \begin{tabularx}{\textwidth}{>{\bfseries}l c X}
        \toprule
        \textbf{Komponen Kerja} & \textbf{Status} & \textbf{Capaian Kunci \& Temuan Empiris} \\
        \midrule
        1. Pengumpulan \& Ground Truth & 100\% & Integrasi 13 GeoJSON Klaten + 7 GeoJSON Salatiga + non-sawah; verifikasi 52 klaster representatif (16 Srinuk, 5 Inpari 32, 1 Membramo, 30 non-sawah). \\
        2. Segmentasi Spasial SNIC     & 100\% & Parameter optimal ($\text{size}=9$, $\text{compactness}=5$) menghasilkan 2.921 klaster superpixel Delanggu yang mematuhi batas pematang sawah. \\
        3. Ekstraksi Deret Waktu GEE   & 100\% & Rekonstruksi 1.404 titik observasi multi-sensor (1988--2026) dan sub-matriks fokus Sentinel-2 2018--2025 (675 titik observasi). \\
        4. Pemodelan KNN-DTW           & 100\% & Akurasi Tahap 1 (Sawah vs Non-Sawah): \textbf{98,08\%}; Akurasi Tahap 2 (Srinuk vs Non-Srinuk): \textbf{94,23\%} (\textit{Recall} Srinuk \textbf{93,75\%}). \\
        5. Visualisasi WebGIS \& Plot  & 100\% & Dashboard WebGIS interaktif (Leaflet + Chart.js) dan 110+ gambar deret waktu 300 DPI lengkap dengan galeri web. \\
        6. Penulisan Naskah Skripsi    & 55\%  & Naskah Bab 1 (Pendahuluan) dan Bab 2 (Landasan Teori) selesai terstruktur dalam template resmi FTIS UNPAR; Bab 3 \& 4 siap disusun. \\
        \bottomrule
    \end{tabularx}
\end{table}

% -----------------------------------------------------------------------------
% 2. REKAPITULASI AKTIVITAS YANG TELAH DISELESAIKAN
% -----------------------------------------------------------------------------
\section{Rekapitulasi Tahapan dan Kegiatan yang Telah Dilakukan}

\subsection{Pengolahan dan Harmonisasi Data Kebenaran Lapangan (Ground Truth)}
\begin{itemize}
    \item Menginventarisasi dan membersihkan 13 berkas GeoJSON survei lapangan resmi di Klaten serta 7 berkas GeoJSON pembanding di Salatiga (Tingkir).
    \item Melakukan standardisasi skema atribut ke dalam satu berkas terpadu \texttt{unified\_ground\_truth\_all.geojson}, mencakup informasi varietas (\textit{Rojolele Srinuk, Inpari 32, Membramo, Mapan}), fase usia tanam (1 minggu s.d. 3 bulan), koordinat spasial metrik UTM Zone 49S (EPSG:32749), dan kelas tutupan (\textit{sawah} vs \textit{non-sawah}).
    \item Melakukan \textit{spatial overlay} terhadap klaster superpixel dengan batas ambang irisan ($\ge 25\%$), menghasilkan 52 klaster terverifikasi lapangan untuk pelatihan dan evaluasi model.
\end{itemize}

\subsection{Optimasi Segmentasi Citra Berbasis Objek (SNIC Superpixel)}
\begin{itemize}
    \item Mengembangkan algoritma pencarian kisi (*grid search*) parameter SNIC pada platform Google Earth Engine (\texttt{01\_snic\_parameter\_optimizer.py}) menggunakan komposit citra Sentinel-2 bebas awan musim kemarau (21 Agustus 2020, tutupan awan 0{,}11\%).
    \item Menggunakan fungsi optimasi multi-objektif Pareto terbobot:
    \begin{equation}
        \text{Score} = \varpi \cdot \overline{\text{Var}}_{\text{norm}} + (1 - \varpi) \cdot K_{\text{norm}}, \quad (\varpi = 0{,}3)
    \end{equation}
    di mana $\overline{\text{Var}}_{\text{norm}}$ adalah variansi EVI intra-klaster dan $K_{\text{norm}}$ adalah jumlah total klaster.
    \item Menghasilkan parameter terbaik pada $\text{size} = 9$ dan $\text{compactness} = 5$, yang membagi Kecamatan Delanggu menjadi 2.921 klaster poligon alami dengan tingkat reduksi komputasi sebesar 97\% dibandingkan pemrosesan piksel mentah.
\end{itemize}

\subsection{Rekonstruksi Deret Waktu Multi-Sensor Bebas Awan (1988--2026)}
\begin{itemize}
    \item Membangun pipeline otomatis ekstraksi EVI (\textit{Enhanced Vegetation Index}) dan NDVI dari tiga generasi sensor satelit:
    \begin{enumerate}
        \item \textbf{Landsat 5 TM} (1988--2012, 30 m) dengan penapisan bitmask \texttt{QA\_PIXEL} dan resampling bikubik 10 m.
        \item \textbf{Landsat 7 ETM+ \& Landsat 8 OLI} (2012--2017, 30 m) dengan harmonisasi skala pantulan permukaan.
        \item \textbf{Sentinel-2 MSI Level-2A} (2018--2026, 10 m) dengan penapisan modern \textbf{Google Cloud Score+} (\texttt{cs} $\ge 0{,}60$).
    \end{enumerate}
    \item Menghasilkan dataset deret waktu historis panjang sebanyak \textbf{1.404 titik observasi} tanggal bebas awan (\texttt{timeseries\_evi\_matrix\_delanggu.csv}).
    \item Mengekstrak sub-matriks khusus era Sentinel-2 modern (2018--2025: \textbf{675 titik observasi}) untuk analisis fenologi bebas distorsi resolusi spasial sensor lawas.
\end{itemize}

\subsection{Pengembangan dan Visualisasi Produk Diseminasi}
\begin{itemize}
    \item \textbf{Aplikasi WebGIS Interaktif (\texttt{peta\_srinuk\_timeseries\_1988\_2026.html})}: Peta berbasis Leaflet dan Chart.js dengan citra satelit resolusi tinggi Google Hybrid, poligon interaktif per klaster, filter era waktu, serta tombol \textbf{''Tampilkan Rata-Rata''} yang menampilkan profil fenologi rata-rata dari 16 klaster Srinuk (total 20,5 Ha).
    \item \textbf{Galeri Grafik Deret Waktu}: Membangun 52 grafik individual 3-panel di direktori \texttt{grafik\_evi\_cluster\_allvarietas/} dan 52 grafik fokus Sentinel-2 di \texttt{grafik\_evi\_2018\_2025/} (300 DPI) lengkap dengan dasbor peramban web \texttt{index.html}.
\end{itemize}

% -----------------------------------------------------------------------------
% 3. PEMBAHASAN MENDALAM DTW DAN KNN
% -----------------------------------------------------------------------------
\section{Pembahasan Teoretis \& Metodologis: DTW dan KNN-DTW}

\subsection{Urgensi DTW: Masalah Asinkronitas Tanam (*Phase Shift*) di Lapangan}
Dalam budidaya padi sawah tropis, petani di petak-petak yang berdampingan jarang sekali melakukan penanaman serentak pada hari yang sama. Keterbatasan traktor pembajak tanah, ketersediaan air irigasi, dan jadwal semai bibit menimbulkan jeda tanam (*phase shift*) berkisar antara 2 hingga 4 minggu. 

Jika digunakan metrik jarak Euclidean konvensional, kurva deret waktu dibandingkan secara kaku pada titik waktu $t_i$ yang sama:
\begin{equation}
    d_{\text{Euclidean}}(X, Y) = \sqrt{\sum_{i=1}^N (x_i - y_i)^2}
\end{equation}
Ketika Petak A menanam 3 minggu lebih awal daripada Petak B (meskipun keduanya sama-sama menanam varietas Srinuk), pada titik waktu $t$, Petak A telah berada di puncak kanopi vegetatif ($\text{EVI} \approx 0{,}75$) sedangkan Petak B masih berada pada fase anakan awal berair ($\text{EVI} \approx 0{,}30$). Jarak Euclidean akan menghasilkan nilai selisih kuadrat yang sangat besar, sehingga model secara keliru menganggap kedua petak tersebut merupakan varietas yang berbeda.

\subsection{Mekanisme Matematis Dynamic Time Warping (DTW)}
\textit{Dynamic Time Warping} (DTW) memecahkan kendala tersebut melalui perataan sumbu waktu non-linear dengan meregangkan (*stretching*) atau memampatkan (*compressing*) sumbu waktu secara elastis. DTW menyelaraskan titik-titik fenologi yang bersesuaian secara biologis (fase genangan ke genangan, dan puncak kanopi ke puncak kanopi).

Diberikan dua deret waktu $X = (x_1, x_2, \dots, x_N)$ dan $Y = (y_1, y_2, \dots, y_M)$:
\begin{enumerate}
    \item Dihitung matriks jarak lokal antar elemen: $d(x_i, y_j) = (x_i - y_j)^2$.
    \item Dibangun matriks akumulasi biaya kumulatif $D(i, j)$ menggunakan pemrograman dinamis:
    \begin{equation}
        D(i, j) = d(x_i, y_j) + \min \begin{cases}
            D(i-1, j)   & \text{(Peregangan waktu / Insertion)} \\
            D(i, j-1)   & \text{(Kompresi waktu / Deletion)} \\
            D(i-1, j-1) & \text{(Pencocokan titik / Match)}
        \end{cases}
    \end{equation}
    dengan kondisi awal $D(1, 1) = d(x_1, y_1)$ dan jalur optimal berakhir pada $D(N, M)$. Jarak akhir DTW adalah biaya kumulatif total: $d_{\text{DTW}}(X, Y) = D(N, M)$.
\end{enumerate}

\subsection{Kendala Global Pita Sakoe-Chiba (*Sakoe-Chiba Band*)}
Tanpa pembatasan, algoritma DTW dapat memetakan titik waktu yang terpaut terlalu jauh dan tidak masuk akal secara biologis (misalnya menyelaraskan bibit pada bulan Januari dengan fase panen pada bulan Juni). Untuk mencegah distorsi tersebut, diterapkan kendala jendela pita Sakoe-Chiba:
\begin{equation}
    |i - j| \le R
\end{equation}
di mana $R$ adalah radius jendela (*window radius*).
\begin{itemize}
    \item \textbf{Kesesuaian Fenologis}: Membatasi elastisitas pergeseran waktu maksimal $R$ timestep (pada penelitian ini digunakan $R = 15$ hingga $R = 30$), yang mencerminkan toleransi pergeseran jadwal tanam maksimal 1 s.d. 1{,}5 bulan.
    \item \textbf{Efisiensi Algoritma}: Memangkas kompleksitas komputasi dari $O(N \cdot M)$ menjadi linear $O(R \cdot (N+M))$, memungkinkan eksekusi klasifikasi ribuan deret waktu secara efisien.
\end{itemize}

\subsection{Integrasi K-Nearest Neighbors Berbasis DTW (KNN-DTW)}
Algoritma KNN merupakan metode pembelajaran berbasis memori (*instance-based learning*). Pada penelitian ini, fungsi jarak standar Euclidean digantikan oleh fungsi jarak elastis DTW berkendala Sakoe-Chiba:
\begin{equation}
    \text{Distance}(X, Y) = d_{\text{DTW-SakoeChiba}}(X, Y)
\end{equation}
Untuk mengklasifikasikan klaster deret waktu baru $X_{\text{uji}}$, model menghitung jarak DTW ke seluruh data latih, memilih $k$-tetangga terdekat dengan jarak minimum, dan menentukan kelas prediksi melalui pemungutan suara mayoritas (*majority voting*).

\subsection{Arsitektur Klasifikasi Berjenjang Dua Tahap (Two-Stage Classifier)}
Klasifikasi multi-kelas langsung (1 tahap) terbukti kurang efektif karena variasi spektral antara lahan sawah vs objek non-sawah jauh lebih dominan dibandingkan variasi halus antarsesama varietas padi. Oleh karena itu, diterapkan dekomposisi hierarkis dua tahap:
\begin{enumerate}
    \item \textbf{Tahap 1 (Sawah vs Non-Sawah)}:
    \begin{itemize}
        \item Memanfaatkan pola osilasi musiman periodik (lembah genangan air $\rightarrow$ fase vegetatif curam $\rightarrow$ puncak kanopi $\rightarrow$ penurunan pematangan). Objek non-sawah (permukiman, jalan, badan air) memiliki kurva EVI datar atau berfluktuasi tanpa periodisitas tanam.
        \item Parameter terbaik: $k=7$ (atau $k=3$) dengan radius $R=30$.
    \end{itemize}
    \item \textbf{Tahap 2 (Rojolele Srinuk vs Varietas Pembanding)}:
    \begin{itemize}
        \item Membedakan profil pertumbuhan varietas lokal unggul Srinuk (umur tanam 115--120 hari, tajuk daun lebat, kurva puncak lebih lebar) terhadap varietas unggul genjah nasional seperti Inpari 32 (umur 105--112 hari, kenaikan vegetatif cepat, puncak kurva sempit).
        \item Parameter terbaik: $k=5$ dengan radius $R=15$ s.d. $R=30$.
    \end{itemize}
\end{enumerate}

% -----------------------------------------------------------------------------
% 4. HASIL KUANTITATIF EVALUASI KINERJA MODEL
% -----------------------------------------------------------------------------
\section{Hasil Evaluasi Kinerja Model Klasifikasi}

Evaluasi dijalankan menggunakan \textit{Stratified K-Fold Cross-Validation} pada 52 klaster terverifikasi lapangan di Kecamatan Delanggu (30 klaster non-sawah, 16 klaster Rojolele Srinuk, 6 klaster varietas pembanding Inpari 32 dan Membramo).

\subsection{Kinerja Klasifikasi Tahap 1 (Sawah vs Non-Sawah)}
\begin{table}[H]
    \centering\small
    \caption{Metrik Kinerja Evaluasi Tahap 1 (Sawah Aktif vs Non-Sawah)}
    \begin{tabularx}{\textwidth}{>{\raggedright\arraybackslash}Xcccc}
        \toprule
        \textbf{Kelas Target} & \textbf{Precision} & \textbf{Recall} & \textbf{F1-Score} & \textbf{Jumlah Sampel} \\
        \midrule
        Non-Sawah (Permukiman, Air, Jalan) & 1,000 & 0,967 & 0,983 & 30 \\
        Lahan Sawah Aktif                  & 0,957 & 1,000 & 0,978 & 22 \\
        \midrule
        \textbf{Akurasi Keseluruhan (Tahap 1)} & \multicolumn{4}{c}{\textbf{98,08\%} (51 dari 52 klaster terdeteksi tepat)} \\
        \textbf{Weighted Average}            & \textbf{0,982} & \textbf{0,981} & \textbf{0,981} & 52 \\
        \bottomrule
    \end{tabularx}
\end{table}

\subsection{Kinerja Klasifikasi Tahap 2 (Rojolele Srinuk vs Varietas Pembanding)}
\begin{table}[H]
    \centering\small
    \caption{Metrik Kinerja Evaluasi Tahap 2 (Rojolele Srinuk vs Non-Srinuk)}
    \begin{tabularx}{\textwidth}{>{\raggedright\arraybackslash}Xcccc}
        \toprule
        \textbf{Kelas Varietas} & \textbf{Precision} & \textbf{Recall} & \textbf{F1-Score} & \textbf{Jumlah Sampel} \\
        \midrule
        Varietas Pembanding (Inpari 32 / Membramo) & 0,833 & 0,833 & 0,833 & 6 \\
        \textbf{Padi Rojolele Srinuk}              & \textbf{0,938} & \textbf{0,938} & \textbf{0,938} & 16 \\
        \midrule
        \textbf{Akurasi Sawah Tahap 2}       & \multicolumn{4}{c}{\textbf{90,91\%} (20 dari 22 klaster sawah terdeteksi tepat)} \\
        \textbf{Akurasi Keseluruhan Pipeline} & \multicolumn{4}{c}{\textbf{94,23\%} (49 dari 52 klaster end-to-end)} \\
        \textbf{Weighted F1-Score Pipeline}  & \multicolumn{4}{c}{\textbf{0,941}} \\
        \bottomrule
    \end{tabularx}
\end{table}

\noindent
\textbf{Temuan Evaluasi}:
\begin{itemize}
    \item Model berhasil mendeteksi \textbf{15 dari 16 klaster Rojolele Srinuk} secara presisi (\textit{Recall} Srinuk = \textbf{93,75\%}).
    \item Satu klaster Srinuk yang mengalami misklasifikasi (\texttt{cluster\_id: 2835}) disebabkan oleh poligon sawah yang berbatasan langsung dengan pemukiman padat dan pohon peneduh, sehingga nilai EVI terdistorsi pantulan spektral kanopi pohon.
\end{itemize}

% -----------------------------------------------------------------------------
% 5. INVENTARIS PRODUK & ARTEFAK
% -----------------------------------------------------------------------------
\section{Inventaris Produk \& Artefak yang Telah Dihasilkan}

\begin{enumerate}
    \item \textbf{Modul Kode Sumber Python \& GEE}:
    \begin{itemize}
        \item \texttt{src/01\_snic\_parameter\_optimizer.py}: Algoritma pencarian parameter optimal SNIC di GEE.
        \item \texttt{src/04\_knndtw\_hierarchical\_classifier.py}: Implementasi model klasifikasi 2-tahap KNN-DTW.
        \item \texttt{scripts/snic\_batch\_extraction\_gee.js}: Skrip batch pemrosesan modular tahunan di GEE.
    \end{itemize}
    \item \textbf{Dataset Olahan Spasial \& Deret Waktu}:
    \begin{itemize}
        \item \texttt{timeseries\_evi\_matrix\_delanggu.csv} (1.404 timestep 1988--2026).
        \item \texttt{timeseries\_evi\_sentinel2\_2018\_2025.csv} (675 timestep Sentinel-2).
        \item \texttt{unified\_ground\_truth\_all.geojson} \& \texttt{delanggu\_snic\_clusters\_labeled.geojson}.
    \end{itemize}
    \item \textbf{Aplikasi WebGIS Interaktif}:
    \begin{itemize}
        \item \texttt{outputs/maps/peta\_srinuk\_timeseries\_1988\_2026.html} (WebGIS Leaflet + Chart.js).
    \end{itemize}
    \item \textbf{Galeri Visual Publikasi}:
    \begin{itemize}
        \item Direktori \texttt{grafik\_evi\_cluster\_allvarietas/}: 52 grafik klaster 3-panel + 5 plot komparatif.
        \item Direktori \texttt{grafik\_evi\_2018\_2025/}: 52 grafik fokus EVI Sentinel-2 + 4 komparasi + dasbor galeri web.
    \end{itemize}
    \item \textbf{Draf Naskah Skripsi LaTeX}:
    \begin{itemize}
        \item \texttt{SKRIPSI HAKIM/Bab/bab1.tex} (Pendahuluan lengkap).
        \item \texttt{SKRIPSI HAKIM/Bab/bab2.tex} (Landasan teori agronomis, SNIC, DTW, KNN, WebGIS).
    \end{itemize}
\end{enumerate}

% -----------------------------------------------------------------------------
% 6. KENDALA & RENCANA TAHAP SELANJUTNYA
% -----------------------------------------------------------------------------
\section{Kendala yang Dihadapi \& Rencana Kerja Selanjutnya}

\subsection{Kendala dan Mitigasi yang Diterapkan}
\begin{enumerate}
    \item \textbf{Keterbatasan Poligon Ground Truth Non-Delanggu}: Data Inpari 32 dari Salatiga memiliki karakteristik elevasi dan agroklimat berbeda dengan Delanggu. \textit{Solusi}: Dilakukan kalibrasi spasial lokal dan integrasi sampel Inpari 32 lokal Klaten sebagai data latih utama.
    \item \textbf{Dinamika Tutupan Awan Tropis}: Tingginya tutupan awan saat musim tanam penghujan. \textit{Solusi}: Diterapkan penapisan Google Cloud Score+ ($\ge 0{,}60$) dipadukan dengan teknik interpolasi temporal dan \textit{gap-filling} bebas nilai kosong (\textit{NaN-free}).
    \item \textbf{Kerapian Visualisasi Grafik}: Terdapat collision/overlap teks judul pada grafik berdimensi besar. \textit{Solusi}: Memodifikasi engine perender Matplotlib dengan pengaturan eksplisit \texttt{gridspec} dan hierarki margin atas.
\end{enumerate}

\subsection{Rencana Kerja Tahap Selanjutnya (Roadmap)}
\begin{table}[H]
    \centering\small
    \caption{Roadmap Penyelesaian Tugas Akhir}
    \begin{tabularx}{\textwidth}{lXcc}
        \toprule
        \textbf{Target Aktivitas} & \textbf{Uraian Pekerjaan} & \textbf{Status} & \textbf{Target Waktu} \\
        \midrule
        \textbf{Bab 3 (Metodologi)} & Penulisan alur penelitian, formula matematis SNIC, DTW, dan KNN & Siap Ditulis & Minggu I Oktober 2026 \\
        \textbf{Bab 4 (Hasil \& Bahas)} & Penyajian tabel metrik, confusion matrix, profil kurva, dan WebGIS & Siap Ditulis & Minggu II Oktober 2026 \\
        \textbf{Bab 5 (Kesimpulan)}   & Perumusan kesimpulan evaluasi model dan saran pengembangan & Siap Ditulis & Minggu III Oktober 2026 \\
        \textbf{Finalisasi Skripsi}   & Penyelarasan daftar pustaka BibTeX, tata letak, dan draf sidang & Terjadwal & Minggu IV Oktober 2026 \\
        \bottomrule
    \end{tabularx}
\end{table}

\vspace{0.4cm}
\noindent
Demikian laporan kemajuan penelitian tugas akhir ini disusun dengan sebenar-benarnya sebagai pertanggungjawaban progres akademik kepada Dosen Pembimbing.

\vspace{0.6cm}
\begin{flushright}
    Bandung, 30 September 2026\\[1.2cm]
    \textbf{Muhammad Hakim Ghifari Akiyat}\\
    NPM: 6182201096
\end{flushright}

\end{document}
